\documentclass[UTF8,a4paper,zihao=-4]{ctexart}

% 2026 第十一届全国密码技术竞赛作品设计报告 LaTeX 模板
% 编译方式：latexmk -xelatex -interaction=nonstopmode -halt-on-error template.tex

\usepackage[top=2cm,bottom=2cm,left=2cm,right=2cm,
            headheight=15pt,headsep=1.5pt,footskip=24pt]{geometry}
\usepackage{fontspec}
\usepackage{xeCJK}
\usepackage{graphicx}
\usepackage{array}
\usepackage{tabularx}
\usepackage{fancyhdr}
\usepackage{titlesec}
\usepackage{enumitem}
\usepackage{amsmath,amssymb}
\usepackage{booktabs}
\usepackage{caption}
\usepackage{listings}
\usepackage{xcolor}
\usepackage{etoolbox}
\usepackage{pifont}
\usepackage{hyperref}

% ---------- 字体 ----------
% 与 Word 模板对应：楷体、华文楷体、黑体、宋体。
\setCJKmainfont[AutoFakeBold=2]{KaiTi}
\setmainfont{Times New Roman}
\newCJKfontfamily\coverfont[AutoFakeBold=2]{STKaiti}
\newCJKfontfamily\heifont[AutoFakeBold=2]{SimHei}
\newCJKfontfamily\songfont{SimSun}

% ---------- 作品信息：仅需修改本区域 ----------
\newcommand{\WorkNumber}{}
\newcommand{\WorkTitle}{}
\newcommand{\WorkDate}{\hspace{1.5em}年\hspace{1.5em}月\hspace{1.5em}日}
% 可选值：软件设计、硬件制作、工程实践、密码应用技术、其它；留空则全部不勾选。
\newcommand{\SelectedCategory}{}
\newcommand{\WorkAbstract}{}
\newcommand{\WorkKeywords}{}

% ---------- 通用格式 ----------
\hypersetup{hidelinks,pdfstartview=FitH}
\setlength{\parindent}{2em}
\setlength{\parskip}{0pt}
% 固定行距不再乘全局倍率；标题的行距单独设置。
\linespread{1}
\AtBeginDocument{\fontsize{14pt}{24pt}\selectfont}
\setlist{nosep,leftmargin=2em}
\captionsetup{font=small,labelsep=quad}

\newcommand{\emptybox}{\makebox[1em][l]{$\square$}}
\newcommand{\checkedbox}{\makebox[1em][l]{%
  \rlap{$\square$}\raisebox{0.05ex}{\ding{51}}}}
\newcommand{\categorybox}[1]{%
  \ifdefstring{\SelectedCategory}{#1}{\checkedbox}{\emptybox}#1}
% 空白模板项之间的占位间距；填写正文时可删除。
\newcommand{\templategap}{\mbox{}\par\vspace{0.1cm}}

% 正文页眉、页脚与 Word 模板一致。
\fancypagestyle{mainstyle}{%
  \fancyhf{}
  \fancyhead[C]{{\songfont\zihao{-5} 全国密码技术竞赛作品设计报告}}
  \fancyfoot[C]{\raisebox{6.7mm}[0pt][0pt]{\fontsize{9pt}{10pt}\selectfont\thepage}}
  \fancyheadoffset[L,R]{0.53mm}
  \renewcommand{\headrulewidth}{0.75pt}
  \renewcommand{\footrulewidth}{0pt}
}

% 一级标题：15 pt 华文楷体；二级标题：14 pt 黑体。
\titleformat{\section}
  {\coverfont\bfseries\fontsize{15pt}{22.5pt}\selectfont}
  {\thesection.}{0pt}{}
\titlespacing*{\section}{0pt}{12pt}{12pt}
\titleformat{\subsection}
  {\heifont\bfseries\fontsize{14pt}{21pt}\selectfont}
  {\thesubsection}{0.7em}{}
\titlespacing*{\subsection}{0pt}{12pt}{12pt}
\setcounter{secnumdepth}{2}

\lstset{
  basicstyle=\small\ttfamily,
  frame=single,
  breaklines=true,
  columns=fullflexible,
  showstringspaces=false
}

% 固定高度的基本信息表单元格。
\newcommand{\formcell}[2]{%
  \parbox[t][#1][t]{\linewidth}{\vspace{-0.35em}#2}}

\begin{document}

% ==================== 封面 ====================
\hypersetup{pageanchor=false}
\thispagestyle{empty}
\noindent
\setlength{\unitlength}{1mm}
\begin{picture}(169,257)
  \put(0,252.9){\makebox(0,0)[l]{\fontsize{14pt}{24pt}\selectfont 作品编号：}}
  \put(23.5,251.6){\makebox(83,0)[l]{\scalebox{1.095}{\fontsize{14pt}{24pt}\selectfont
    \ifdefempty{\WorkNumber}{\color{red}\bfseries（必须填写系统自动分配的编号）}{\color{black}\WorkNumber}}}}
  \put(24,249.6){\line(1,0){81.8}}

  \put(0,234){\makebox(0,0)[l]{\fontsize{14pt}{24pt}\selectfont 作品类别：}}
  \put(27,234){\makebox(0,0)[l]{\songfont\fontsize{10.5pt}{24pt}\selectfont
    \categorybox{软件设计}\quad\categorybox{硬件制作}\quad\categorybox{工程实践}}}
  \put(27,225.565){\makebox(0,0)[l]{\songfont\fontsize{10.5pt}{24pt}\selectfont
    \categorybox{密码应用技术}\quad\categorybox{其它}}}

  \put(126,232){\includegraphics[width=42.7mm,height=32.3mm,keepaspectratio]{imgs/logo-2026.jpeg}}

  \put(15,169){\makebox(140,0)[c]{\bfseries\fontsize{18pt}{24pt}\selectfont
    2026年第十一届全国密码技术竞赛作品设计报告}}
  \put(15,158){\rule{140mm}{1.2pt}}

  \put(30,132){\makebox(0,0)[l]{\fontsize{14pt}{25pt}\selectfont 作品题目：}}
  \put(61,130.5){\line(1,0){90}}
  \put(62,132){\makebox(88,0)[c]{\fontsize{14pt}{25pt}\selectfont \WorkTitle}}

  \put(0,108){\parbox[t]{166mm}{\bfseries\fontsize{14pt}{24pt}\selectfont
    \hspace{2em}注意：作品为匿名评审，作品设计报告封面和内容中严禁出现单位名称、
    指导老师姓名或团队成员姓名，违者将被取消参赛资格。提交作品时请将此段落文字删除。}}

  \put(0,45){\makebox(170,0)[c]{\fontsize{14pt}{24pt}\selectfont \WorkDate}}
  \put(0,34){\makebox(170,0)[c]{\fontsize{14pt}{24pt}\selectfont 中国密码学会}}
\end{picture}

% ==================== 基本信息表 ====================
\clearpage
\hypersetup{pageanchor=true}
\pagestyle{mainstyle}
\setcounter{page}{1}
\thispagestyle{mainstyle}
\vspace*{-24.6pt}
\begin{center}
\songfont
\renewcommand{\arraystretch}{1}
\setlength{\tabcolsep}{2mm}
\setlength{\arrayrulewidth}{0.5pt}
\begin{tabular}{!{\vrule width 1.5pt}p{15.22cm}!{\vrule width 1.5pt}}
\noalign{\hrule height 1.5pt}
\formcell{0.55cm}{\centering\fontsize{15pt}{22.5pt}\selectfont 基本信息表}\\ \hline
\formcell{0.57cm}{\fontsize{15pt}{22.5pt}\selectfont 编号：\WorkNumber}\\ \hline
\formcell{0.57cm}{\fontsize{15pt}{22.5pt}\selectfont 作品题目：\WorkTitle}\\ \hline
\formcell{0.48cm}{\fontsize{15pt}{22.5pt}\selectfont 作品类别：\quad
  \fontsize{10.5pt}{18pt}\selectfont
  \categorybox{软件设计}\quad\categorybox{硬件制作}\quad\categorybox{工程实践}\quad
  \categorybox{密码应用技术}\quad\categorybox{其它}}\\ \hline
\formcell{16.23cm}{\fontsize{15pt}{22.5pt}\selectfont 作品内容摘要：\par\vspace{0.5em}
  {\fontsize{10.5pt}{12.6pt}\selectfont\setlength{\parindent}{2em}%
   \ifdefempty{\WorkAbstract}{}{\indent\WorkAbstract\par}}}\\ \hline
\formcell{2.20cm}{\fontsize{15pt}{22.5pt}\selectfont 关键词（五个）：\par\vspace{0.35em}
  {\fontsize{10.5pt}{12.6pt}\selectfont \WorkKeywords\par}}\\
\noalign{\hrule height 1.5pt}
\end{tabular}
\end{center}

% ==================== 正文 ====================
\clearpage
% 调整空白模板的正文页首位置。
\vspace*{-9mm}
\section{作品功能与性能说明}
% 在此填写作品功能、性能和创新点。
\templategap

\section{设计与实现方案}
\templategap

\subsection{实现原理}
（硬件框图、软件流程、相关描述等）
\mbox{}\par\vspace{0.7cm}

\subsection{运行结果}
\templategap

\subsection{技术指标}
\templategap

\section{系统测试与结果}
\templategap

\subsection{测试方案}
\templategap

\subsection{功能测试}
\templategap

% 官方模板在此处换页。
\clearpage
\subsection{性能测试}
\templategap

\subsection{测试数据与结果}
\templategap

\section{应用前景}
\templategap

\section{结论}
\vspace{1.9cm}

{\fontsize{12pt}{24pt}\selectfont
注意：这里列出的一级标题和二级标题供参考，参赛团队可以根据作品实际情况适当增删一级标题和二级标题。}

\end{document}
